\documentclass[english]{article}
\usepackage[T1]{fontenc}
\usepackage[utf8]{inputenc}
\usepackage{babel}
\usepackage[unicode=true,pdfusetitle,
 bookmarks=true,bookmarksnumbered=false,bookmarksopen=false,
 breaklinks=true,pdfborder={0 0 1},backref=false,colorlinks=false]
 {hyperref}
\usepackage{tabularx}
\usepackage{graphicx}
\graphicspath{{images/}}
\usepackage{svg}
\usepackage{float}
\usepackage{titling}
\renewcommand{\arraystretch}{1.4}
\newcommand{\code}[1]{\texttt{#1}}
\usepackage{listings}
\usepackage{color}
\definecolor{javared}{rgb}{0.6,0,0} % for strings
\definecolor{javagreen}{rgb}{0.25,0.5,0.35} % comments
\definecolor{javapurple}{rgb}{0.5,0,0.35} % keywords
\definecolor{javadocblue}{rgb}{0.25,0.35,0.75} % javadoc
\definecolor{cornellred}{rgb}{0.7, 0.11, 0.11}

\lstset{language=Java,
	basicstyle=\ttfamily\small,
	xleftmargin=0cm,
	keywordstyle=\color{javapurple}\bfseries,
	stringstyle=\color{javared},
	commentstyle=\color{javagreen},
	morecomment=[s][\color{javadocblue}]{/**}{*/},
	numbers=left,
	numberstyle=\tiny\color{black},
	stepnumber=1,
	numbersep=10pt,
	tabsize=4,
	showspaces=false,
	showstringspaces=false,
	linewidth=13cm,
	breaklines=true
	}

\pretitle{%
	\begin{center}
		\LARGE
		\includegraphics[width=250pt]{../other/Logo_blu.png}\\[\bigskipamount]~\\[\bigskipamount]
	}
\posttitle{\end{center}}

\begin{document}

\title{Politecnico di Milano\\
 A.A. 2016–2017 \\
Software Engineering 2: “PowerEnJoy” \\
\emph{\textbf{Code Inspection Document}}}

\author{Pietro Ferretti, Nicole Gervasoni, Danilo Labanca}
\date{February 5, 2017}
\maketitle

\newpage

\tableofcontents{}

\newpage

\section{Assigned Class}
% <state the namespace pattern and name of the classes that were assigned to you.>
We were assigned one class to analyse, located in:\\

\texttt{apache-ofbiz-16.11.01/framework/service/\\
\hspace*{1cm}src/main/java/org/apache/ofbiz/service/job/JobManager.java}

\section{Functional Role}% of assigned set of classes: <elaborate on the functional role you have identified for the class cluster that was assigned to you, also, elaborate on how you managed to understand this role and provide the necessary evidence, e.g., javadoc, diagrams, etc.>

We found some information about the class we've been assigned and the environment it belongs to on Apache OFBiz's website\footnote{\href{https://ofbiz.apache.org/}{https://ofbiz.apache.org/}}
and its wiki\footnote{\href{https://cwiki.apache.org/confluence/display/OFBIZ/Home}{https://cwiki.apache.org/confluence/display/OFBIZ/Home}}.

\subsection{Apache OFBiz}
\paragraph{}
Apache OfBiz is an open source ERP system for the automation of enterprise processes. It offers many different applications and components to integrate and manage business processes like resource management, activities planning and customer relationship management.


\subsection{Service Engine Framework}
\paragraph{}
The Service Engine is one of the components of Apache OFBiz.

\paragraph{}
Services are independent pieces of logic which when placed together process many different types of business requirements. Services can be of many different types: Workflow, Rules, Java, SOAP, BeanShell, etc.

\paragraph{}
Services are defined through the Service Definition and are assigned to a specific Service Engine. Each Service Engine is responsible for invoking the defined service in an appropriate way. Services can be scheduled to run at specific times to run in the background via the Job Scheduler.

\subsection{Job Scheduler}

\paragraph{}
The Job Scheduler (also called Job Manager) is integrated with the services framework, and is tasked with accepting new jobs to schedule and running them at the correct time.

\paragraph{}
The scheduler is a multi-threaded component with a single thread used for job managing/scheduling and separate threads used for invocation of each service. When a job is scheduled to run, the scheduler will call the service dispatcher associated with the job to invoke the service in its own thread. This will prevent long or time consuming jobs from slowing down other jobs in the queue.


\subsection{Assigned Class: Job Manager}

\paragraph{}
We can find more specific details about the class in the class Javadoc:

\begin{lstlisting}

/**
* Job manager. The job manager queues and manages jobs. Client
* code can queue a job to be run immediately by calling the 
* runJob({@link #runJob(Job)}) method, or schedule a job to be 
* run later by calling the {@link #schedule(String, String,
* String, Map, long, int, int, int, long, int)} method.
* Scheduled jobs are persisted in the JobSandbox entity.
* 
* <p>A scheduled job's start time is an approximation - the
* actual start time will depend on the job manager/job poller
* configuration (poll interval) and the load on the server.
* Scheduled jobs might be rescheduled if the server is busy.
* Therefore, applications requiring a precise job start time 
* should use a different mechanism to schedule the job.</p>
*/
\end{lstlisting}

\paragraph{}
The Javadoc describes two public methods: \texttt{runJob} and \texttt{schedule}, respectively used to execute a job immediately (or at least as soon as the resources are available), and to schedule a job to be run at a later time. 

\paragraph{}
There are two relevant methods in the class that are not mentioned in the main Javadoc:
\begin{itemize}
	\item \texttt{reloadCrashedJobs} is a public method that checks the status of all scheduled jobs in the queue, and reschedules every job that has crashed, if there are any;
	\item \texttt{poll} is a protected method that checks the job queue and returns a list of all the jobs that are scheduled to run.
\end{itemize}

\newpage
\section{List of Issues} % found by applying the checklist: <report the classes/code fragments that do not fulfill some points in the check list. Explain which point is not fulfilled and why.>
% what are we gonna do
\paragraph{}
We list here all the issues we have found in the code following the checklist that we where provided.

\paragraph{}
Points are labeled as "Not ok" if issues are present, and "Ok" if everything is fine.


\subsection{Naming Conventions}

\paragraph{1. Meaningful Names}
Ok. All names are meaningful.

\paragraph{2. One-character variables}
Ok, there are no one-character variables.

\paragraph{3. Class names}
Ok, every class name is in mixed case and properly capitalized.

\paragraph{4. Interface names}
Ok. No interfaces are declared in this file.

\paragraph{5. Method names}
Ok. Every method name is a verb and every method name is camelCase and properly capitalized.

\paragraph{6. Class variables}
Ok. Every class variable is in mixed case and properly capitalized.

\paragraph{7. Constants}
\textcolor{cornellred}{\textbf{Not ok.}}
\paragraph{}
The field 
\begin{lstlisting} 
private static final module 
\end{lstlisting} 
and 
\begin{lstlisting} 
private static final instanceId 
\end{lstlisting} 
are immutable, so they can be considered as constant; they should be capitalized. 

\paragraph{}
On the other hand
\begin{lstlisting} 
private static final registeredManagers 
\end{lstlisting}
 is fine because it's a mutable HashMap.

%https://stackoverflow.com/a/30789385

\subsection{Indentation}

\paragraph{8. Number of spaces}
Ok, the code is consistently indented with 4 spaces.

\paragraph{9. No tabs for indentation}
Ok. No tabs are used to indent the code.

\subsection{Braces}

\paragraph{10. Consistent bracing style}
Ok.
The code is consistently braced following the \textit{"Kernighan and Ritchie"} style.

\paragraph{11. One-line statements bracing}
\textcolor{cornellred}{\textbf{Not ok.}}
\paragraph{}
In the code we can find the following statements that don't follow this rule:
\begin{itemize}
	\item 
	\begin{lstlisting} 
	if (Debug.infoOn()) Debug.logInfo("Scheduling Job : " + job, module); \end{lstlisting}
	 on line 326;
	\item 
	\begin{lstlisting} 
	if (Debug.infoOn())
    Debug.logInfo("-- " + rescheduled + " jobs re-scheduled", module); \end{lstlisting} 
    on line 351-352;
    \item 
    \begin{lstlisting} 
    if (Debug.infoOn())
    Debug.logInfo("No crashed jobs to re-schedule", module); \end{lstlisting}
     on line 354-355;
\end{itemize}

\subsection{File Organization}
\paragraph{12. Blank lines as separation}
Ok.
Blank lines are present between each method, around imports and variable declarations.
Most of the methods also begin with a Javadoc.

\paragraph{13. Where practical, line length under 80 characters}
\textcolor{cornellred}{\textbf{Not ok.}}
\paragraph{}
There is a huge number of lines that exceed 80 characters.

\paragraph{14. Line length always under 120 characters}
\textcolor{cornellred}{\textbf{Not ok.}}
\paragraph{}
The lines \textit{74}, \textit{186}, \textit{198}, \textit{217}, \textit{221}, \textit{222}, \textit{261-264}, \textit{273}, \textit{311}, \textit{315}, \textit{317}, \textit{387}, \textit{409}, \textit{429}, \textit{453}, \textit{498}, \textit{543}, \textit{560}, \textit{561} don't follow this rule. The most long line is the number 221 with 234 characters:

\begin{lstlisting}
List<EntityExpr> updateExpression = UtilMisc.toList(EntityCondition.makeCondition("jobId", EntityOperator.EQUALS, jobValue.get("jobId")), EntityCondition.makeCondition("runByInstanceId", EntityOperator.EQUALS, null));
\end{lstlisting}

\subsection{Wrapping Lines}
\paragraph{15. Line breaks after commas and operators}
\textcolor{cornellred}{\textbf{Not ok.}}
\paragraph{}
Line \textit{152} starts with a stray comma that should be moved to the previous line.

\paragraph{16. Higher-level breaks are used}
Ok. There are no operators that need line breaks.

\paragraph{17. Statements are aligned to previous ones}
Ok. All statements are aligned to previous ones.

\subsection{Comments}

\paragraph{18. Comments use}
\textcolor{cornellred}{\textbf{Not ok.}}
\paragraph{}
The method 
\begin{lstlisting} 
	public synchronized void reloadCrashedJobs() 
\end{lstlisting}

 on line \textit{305} is not commented and so it isn't easy to understand. 
\paragraph{}
The rest of the code is well commented.

\paragraph{19. Commented out code}
Ok. There are no commented out lines of code in the source code.

\subsection{Java Source Files}
\paragraph{20. Single public class or interface}
Ok.
Job manager is the only public class declared in the file.
There are no other classes.

\paragraph{21. The public class is the first class in the file}
Ok.
Job manager is the only public class declared in the file.
There are no other classes.

\paragraph{22. External program interfaces are consistent with the Javadoc}
\textit{Almost ok.}
The description in the first Javadoc is suitable and consistent with the runJob and schedule interfaces.
\paragraph{}
It would be a good idea to also add the public and protected reloadCrashedJobs and poll  methods to the Javadoc.

\newpage
\paragraph{23. The Javadoc is complete}
\textcolor{cornellred}{\textbf{Not ok.}}
\paragraph{}
The javadoc is not complete:
\begin{itemize}
	\item There is no javadoc for 'module' line 71
	\item There is no javadoc for 'instanceId' line 71
	\item There is no javadoc for reloadCrashedJob line 304
	\item Missing @return tag on getInstance, line 88
	\item Missing @return tag on getDelegator, line 119
	\item Missing @return tag on getDispatcher, line 124
	\item Missing @param tag for 'limit' on poll, line 174
	\item Missing @return tag on poll, line 174
	\item Missing @param tag for 'job' on runJob, line 363
	\item Missing @throws tag for 'JobManagerException' on runJob, line 363
	\item Missing @throws tag for 'JobManagerException' on schedule, line 386, 408, 428, 453, 469, 498, 543
\end{itemize}

\subsection{Package and Import Statements}
\paragraph{24. Package statements are first, import statements second}
Ok.
There is one package statement and all import statements immediately follow.

\subsection{Class and Interface Declarations}
\paragraph{25. The class declarations should follow a specific order}
\textcolor{cornellred}{\textbf{Not ok.}}
\paragraph{}

The static method 
\begin{lstlisting}
	private static void assertIsRunning() {
	        if (isShutDown) {
	            throw new IllegalStateException("OFBiz shutting down");
	        }
	    }
\end{lstlisting}
should be moved under the method

\begin{lstlisting}
	public static JobManager getInstance(Delegator delegator, boolean enablePoller)
\end{lstlisting}
that works as the class constructor.

\paragraph{26. Methods are grouped by functionality}
Ok.
Methods can be roughly grouped as follows:
\begin{itemize}
	\item Check status and constructors:
	\begin{itemize}
		\item{assertIsRunning}
		\item{getInstance}
		\item{shutDown}
	\end{itemize}

	\item Getters:
	\begin{itemize}
		\item{getDelegator}
		\item{getDispatcher}
		\item{getPoolState}
	\end{itemize}

	\item Jobs operations
	\begin{itemize}
		\item{isAvailable}
		\item{getRunPools}
		\item{pool}
		\item{reloadCrashedJobs}
		\item{runJob}
		\item{schedule}
	\end{itemize}
\end{itemize}

\paragraph{27. The code is free of duplicates, long methods, big classes, breaking encapsulation, and coupling and cohesion are adequate}
Ok.
The code is free of duplicates, is organised in small classes and encapsulation is not broken.\\
Moreover the methods are short. Some methods conduct the role of delegator or dispatcher and guarantee a low coupling.\\
The cohesion is respected because all the methods are used to run or queue jobs.

\subsection{Initialization and Declarations}

\paragraph{28. Visibility} \textit{Almost Ok.}
All variables and class members are of the correct type and have the proper visibility.
In line 305 the method 
\begin{lstlisting} 
	public synchronized void reloadCrashedJobs()
\end{lstlisting} 
could be stated as protected.

\paragraph{29. Proper scope.} 
Ok. All variables are declared in the proper scope.

\paragraph{30. New objects.} 
Ok. Each time a new object is desired the proper constructor is called.

\paragraph{31. All object references are initialized before use.} 
Ok. All reference are initialized before any object uses it. 

\paragraph{32. Variables initialization.} 
Ok. All variables are initialized where they are declared, unless dependent upon a computation.

\paragraph{33. Declarations.} Ok. Each declaration appear at the beginning of blocks.

\subsection{Method Calls}

\paragraph{34. Correct orders parameters}

Ok. All the methods are called with the parameters in the correct order.

\paragraph{35. The called method is the right method}

Ok. The methods called are always the right ones.

\paragraph{36. The returned value from the method is used properly}

Ok. The value returned from the method is always used accurately.


\subsection{Arrays}
\paragraph{37. No off-by-one errors in array indexing}
Ok.
The only indexing is made with foreach, so there can't be off-by-one errors.

\paragraph{38. No out-of-bounds indexes}
Ok.
There is no number indexing.

\paragraph{39. Constructors are called when a new array item is desired}
Ok.
There are no arrays, only collections. Every collection is created with the appropriate constructor.

\subsection{Object Comparison}
\paragraph{40. Objects are compared with equals}
Ok.
There are no object comparisons.

\newpage
\subsection{Output Format}
\paragraph{41. Displayed output is free of spelling and grammatical errors}
\textcolor{cornellred}{\textbf{Not ok.}}
\paragraph{}
We found two minor mistakes:

\begin{itemize}
	\item{Lines 156: 
	\begin{lstlisting}
		Debug.logWarning(e, "Exception thrown while check lock on JobManager : " + instanceId, module);
	\end{lstlisting}
	it should be "while checking".}
	\item{Lines 156: 
	\begin{lstlisting}
		Debug.logWarning("Unable to locate DispatchContext object; not running job!", module);
	\end{lstlisting}
	it should be "job:", as in the other debug log.}
\end{itemize}


\paragraph{42. Error messages are comprehensive and useful}
Ok. All error messages clearly explain what type of problem has occurred

\paragraph{43. Output is formatted correctly in terms of line breaks and spacing} \textit{Almost Ok.}
There are no outputs that need line breaks but some debug outputs don't have a trailing space.

\subsection{Computation, Comparisons and Assignments}
\paragraph{44."Brutish programming".}
Ok. The implementation avoids brute force solutions; the code is simple and concise.

\paragraph{45. Operator precedence and parenthesizing.}
Ok. Computation/evaluation of operator precedence and parentheses is in the proper order.
 
\paragraph{46. The liberal use of parenthesis is used to avoid operator precedence problems.}
Ok. There isn't any parenthesis used in an inappropriate way.

\paragraph{47. All denominators of a division are prevented from being zero.}
Ok. There are no divisions.

\paragraph{48. Integer arithmetic, especially division, are used appropriately to avoid causing unexpected truncation/rounding.}
Ok. Integer arithmetic is used only to increment variables.

\paragraph{49. Comparison and Boolean operators are correct.}
Ok.

\paragraph{50. Throw-catch expressions.}
Ok. The error conditions are always legitimate.

\paragraph{51. The code is free of any implicit type conversions.}
Ok. The code is free of any implicit type conversion.
\subsection{Exceptions}

\paragraph{52. Relevant exceptions are caught.}
Ok.
\paragraph{53. The appropriate action is taken for each catch block.}
Ok. There are two a little too general
\begin{lstlisting}
catch (Throwable t)
\end{lstlisting} statements, but this is justified as necessary to guarantee a working jobPoller even when a database connection is not available.

\subsection{Flow of Control}
\paragraph{54. All switch cases are addressed with a break}
Ok, no switch statements.

\paragraph{55. All switch statements have a default branch}
Ok, no switch statements.

\paragraph{56. All loops are correctly formed, with appropriate initialization, increments and termination expressions}
Ok.
All for loops are foreach, so no checks on bounds are needed

\subsection{Files}
\paragraph{}
Ok.
The JobManager class does not handle files.

\newpage
\section{Other Problems} % you have highlighted: <list here all the parts
%of code that you think create or may create a bug and explain why.>
The assigned class does not contain any other significant problem.\\
For the sake of accuracy, we suggest to replace in line 576
\begin{lstlisting}
jFields.put("currentRetryCount", new Long(0));\end{lstlisting}
with the more efficient
\begin{lstlisting}
 jFields.put("currentRetryCount", valueOf(0));
\end{lstlisting}
Using \lstinline|new Long()| is guaranteed to always result in a new object whereas \lstinline|Integer.valueOf()| allows caching of values to be done by the compiler, class library, or JVM. Using of cached values avoids object allocation and the code will be faster.\\
This might have to be considered in case of multiple new Long() calls.

\section{Effort Spent}
\begin{itemize}
	\item{Pietro Ferretti: 7 hours of work}
	\item{Nicole Gervasoni: 5 hours of work}
	\item{Danilo Labanca: 5 hours of work}
\end{itemize}


\section{Revisions}

\subsection{Changelog}
\begin{itemize}
	\item{CID v1.0, published on February 5, 2017}
\end{itemize}
\end{document}
